满足严格的首 Token 时间（Time-To-First-Token, TTFT）要求，对 LLM 应用至关重要。为了提升效率，现代 LLM 服务系统采用带有多样化并行方式的解耦架构，引入了包含可复用 KV 块检索、集合通信以及 Prefill-to-Decode（P2D）传输在内的复杂多阶段工作流。不同请求中具有依赖关系的各阶段流会在共享瓶颈链路上发生跨阶段、跨请求重叠，使 TTFT 极易受到网络争用影响，因此亟需具备阶段感知能力的调度机制。不幸的是，现有大多数工作都以与阶段无关的方式调度流，导致缺乏协调的争用，进而成为 SLO 违约的主要原因之一。

本文提出 \sys，一种旨在最大化 TTFT SLO 达成率的整体式\underline{多}阶段\underline{流}\underline{调度}机制。其核心是在无需精确知道请求剩余松弛度的前提下，近似实现 Least-Laxity-First（LLF）调度策略。为此，\sys 通过反向多级队列（Reverse Multi-Level Queue, RMLQ）结构实现了一种 \emph{延后并提升}（Defer-and-Promote）原则。随着有效松弛度不断缩小，\sys 动态提升任务优先级，从而优先保障松弛度更小的流，同时避免 SLO 较宽松的请求过早占用网络带宽。我们将 \sys 实现为可插拔模块，并将其集成到 vLLM 中，在一个由 8 台服务器、32 块 GPU 组成的测试平台以及大规模仿真环境上进行了评估。结果表明，\sys 能够稳定超越当前最先进的基线方法，将 TTFT SLO 达成率提升 1.2$\times$--2.4$\times$。
